\chapter{Befehlsreferenz}
\label{kap:referenz}

Dieser Anhang listet vollständig, was esim versteht: die Programme mit
allen Optionen, jeden Schlüssel der Eingabedatei mit erlaubten Werten und
Standardwerten, alle Ausdrücke und alle Befehle des Editors. Die
Erklärungen stehen in den Kapiteln, auf die verwiesen wird; hier geht es
ums Nachschlagen.

\paragraph{Lesart der Tabellen.}
»Pflicht« heißt, dass der Schlüssel angegeben werden muss. Ein Standardwert
gilt, wenn der Schlüssel fehlt. Wo zwei Namen durch »|« getrennt sind, sind
beide gleichwertig (Synonyme). »Zahl« bedeutet: JSON-Zahl, Zeichenkette mit
SI-Präfix oder Ausdruck aus Parametern (Abschnitt~\ref{sec:zahlen});
»Ausdruck« bedeutet: Zeichenkette mit einem Ausdruck nach
Abschnitt~\ref{sec:ausdruecke}.

% ---------------------------------------------------------------------------
\section{Programme}

\subsection{esim}

\begin{shell}
esim <datei.json> [-o <ordner>] [--summary <datei.json>] [-j <n>] [--progress]
esim -h | --help
\end{shell}

\begin{longtable}{@{}P{4.2cm}P{2.6cm}P{7.0cm}@{}}
  \toprule
  Option & Standard & Wirkung \\
  \midrule
  \endhead
  \bottomrule
  \endfoot
  \code{<datei.json>} & Pflicht & Eingabedatei \\
  \code{-o <ordner>} & aktuelles Verzeichnis & Ordner der Ergebnisdateien \\
  \code{--summary <datei>} & keine & Zusammenfassung aller Ergebnisse als JSON (Abschnitt~\ref{sec:summary}) \\
  \code{-j <n>} | \code{--threads <n>} & alle Kerne & gleichzeitige Läufe einer Studie \\
  \code{--progress} & aus & Zeilen \code{progress <analyse> <anteil>} auf stderr \\
  \code{-h} | \code{--help} & & kurze Hilfe \\
\end{longtable}

Rückgabewert 0 bei Erfolg, 1 bei einem Fehler. Ergebnisdateien:
\datei{<ordner>/<dateiname>.<analyse>.csv} (Tabelle in
Abschnitt~\ref{sec:ausgaenge}ff.).

\subsection{esim-editor}

\begin{shell}
esim-editor [datei.json] [--subcircuit <name>] [--run] [--snapshot <bild.png>]
\end{shell}

\begin{longtable}{@{}P{4.2cm}P{9.6cm}@{}}
  \toprule
  Option & Wirkung \\
  \midrule
  \endhead
  \bottomrule
  \endfoot
  \code{datei.json} & Datei öffnen \\
  \code{--subcircuit <name>} & den Plan dieser Teilschaltung zeigen \\
  \code{--run} & gleich simulieren \\
  \code{--snapshot <bild.png>} & Fenster (nach der Simulation) als Bild sichern und beenden \\
\end{longtable}

\subsection{tools/spice2json.py}

\begin{shell}
python3 tools/spice2json.py <netzliste.cir> [-o <datei.json>]
\end{shell}

\begin{longtable}{@{}P{4.2cm}P{9.6cm}@{}}
  \toprule
  Argument & Wirkung \\
  \midrule
  \endhead
  \bottomrule
  \endfoot
  \code{<netzliste>} & SPICE-Netzliste (\datei{.cir}, \datei{.sp}, \datei{.net}) \\
  \code{-o} | \code{--output <datei>} & JSON-Datei; ohne Angabe die Standardausgabe \\
\end{longtable}

\subsection{tools/plot.py}

\begin{shell}
python3 tools/plot.py <datei.csv> [<datei.csv> ...] [-o <bild>] [--cols <regex>]
                      [--separate] [--group] [--logx] [--logy]
\end{shell}

\begin{longtable}{@{}P{4.2cm}P{9.6cm}@{}}
  \toprule
  Argument & Wirkung \\
  \midrule
  \endhead
  \bottomrule
  \endfoot
  \code{<datei.csv> ...} & eine oder mehrere Ergebnisdateien (übereinandergelegt) \\
  \code{-o} | \code{--output <bild>} & Bild schreiben (PNG, PDF, SVG) statt ein Fenster zu öffnen \\
  \code{--cols <regex>} & nur Spalten, deren Name auf den regulären Ausdruck passt \\
  \code{--separate} & ein Feld je Größe \\
  \code{--group} & erste Spalte als Parameter eines Sweeps behandeln \\
  \code{--logx}, \code{--logy} & logarithmische Achsen \\
\end{longtable}

\subsection{Python-Modul esim}

\begin{longtable}{@{}P{5.2cm}P{8.6cm}@{}}
  \toprule
  Name & Bedeutung \\
  \midrule
  \endhead
  \bottomrule
  \endfoot
  \code{run(netlist, ...)} & simuliert, liefert \code{Result}; weitere Argumente \code{parameters}, \code{analyses}, \code{outputs}, \code{esim}, \code{keep}, \code{threads} (Kapitel~\ref{kap:python}) \\
  \code{load_netlist(pfad)} & Eingabedatei mit Kommentaren als \code{dict} \\
  \code{EsimError} & Ausnahme mit der Fehlermeldung des Simulators \\
  \code{Result.analyses}, \code{.op}, \code{.dc}, \code{.ac}, \code{.tran}, \code{.pss}, \code{.noise}, \code{.sens}, \code{[name]}, \code{.log}, \code{.title} & Analysen und Protokoll \\
  \code{Analysis.type}, \code{.name}, \code{.data}, \code{[spalte]}, \code{.columns}, \code{.values}, \code{.measures}, \code{.stats}, \code{.sweep}, \code{.statistics}, \code{.parameters}, \code{.sensitivities}, \code{.fourier}, \code{.info}, \code{.runs()} & Ergebnisse einer Analyse \\
  Umgebungsvariable \code{ESIM_BIN} & Pfad des Programms \code{esim} \\
\end{longtable}

% ---------------------------------------------------------------------------
\section{Oberste Ebene der Eingabedatei}

\begin{longtable}{@{}P{3.0cm}P{4.3cm}P{6.5cm}@{}}
  \toprule
  Schlüssel & Wert & Bedeutung \\
  \midrule
  \endhead
  \bottomrule
  \endfoot
  \code{title} & Text & Titel \\
  \code{parameters} & \code{{"name": Zahl, ...}} & Parameter (Abschnitt~\ref{sec:parameter}) \\
  \code{functions} & \code{{"f(a, b)": "Ausdruck", ...}} & eigene Funktionen \\
  \code{tables} & \code{{"name": {...}, ...}} & Kennlinien (Abschnitt~\ref{sec:ref-tabellen}) \\
  \code{models} & \code{{"name": Objekt | SPICE-Zeile | Liste von Zeilen}} & Modelle (Abschnitt~\ref{sec:ref-modelle}) \\
  \code{subcircuits} & \code{{"name": Objekt | SPICE-Text | Liste von Zeilen}} & Teilschaltungen (Abschnitt~\ref{sec:ref-teil}) \\
  \code{components} & Liste von Bauteilen & Abschnitt~\ref{sec:ref-bauteile} \\
  \code{variables} & \code{{"name": ...}} & Unbekannte (Abschnitt~\ref{sec:ref-gleichungen}) \\
  \code{equations} & Liste von Gleichungen und Schleifen & eigene Gleichungen \\
  \code{events} & Liste von Ereignissen & Zuweisungen bei Ereignissen \\
  \code{analyses} & Liste von Analysen & Abschnitt~\ref{sec:ref-analysen} \\
  \code{outputs} & Liste von Ausdrücken & Ausgänge; ohne Angabe alle Knoten, Zweigströme, Variablen \\
  \code{options} & Objekt & Abschnitt~\ref{sec:ref-optionen} \\
  \code{nodeset} & \code{{"knoten": Zahl, ...}} & Startwerte für den Arbeitspunkt \\
  \code{ic} & \code{{"knoten": Zahl, ...}} & Spannungen bei $t=0$ für \code{"uic": true} \\
  \code{schematic} & Objekt & Schaltplan des Editors; vom Simulator überlesen \\
\end{longtable}

% ---------------------------------------------------------------------------
\section{Bauteile}
\label{sec:ref-bauteile}

Gemeinsame Schlüssel aller Bauteile:

\begin{longtable}{@{}P{3.0cm}P{4.3cm}P{6.5cm}@{}}
  \toprule
  Schlüssel & Wert & Bedeutung \\
  \midrule
  \endhead
  \bottomrule
  \endfoot
  \code{type} & Text, Pflicht & Bauteiltyp (unten), Groß/klein egal \\
  \code{name} & Text & Name; ohne Angabe Typbuchstabe und laufende Nummer \\
  \code{nodes} & Liste von Knotennamen, Pflicht (außer \code{K}) & Anschlüsse in der Reihenfolge der Tabellen \\
  \code{model} & Text & Modell aus \code{models} (D, Q, J, M) \\
\end{longtable}

\subsection*{Passive Bauteile und Quellen}

\begin{longtable}{@{}P{2.7cm}P{3.0cm}P{3.6cm}P{4.5cm}@{}}
  \toprule
  \code{type} & \code{nodes} & Schlüssel & Standard / Bemerkung \\
  \midrule
  \endhead
  \bottomrule
  \endfoot
  \code{R} | \code{resistor} & a, b & \code{value} | \code{r} & Pflicht \\
   & & \code{noise} & \code{true}; \code{false}: rauschfrei \\
  \code{C} | \code{capacitor} & a, b & \code{value} | \code{c} & Pflicht \\
   & & \code{ic} & Anfangsspannung (nur \code{uic}) \\
  \code{L} | \code{inductor} & a, b & \code{value} | \code{l} & Pflicht \\
   & & \code{ic} & Anfangsstrom (nur \code{uic}) \\
  \code{V} | \code{vsource} & $+$, $-$ & \code{dc} | \code{value} & Wert von \code{wave} bei $t=0$, sonst 0 \\
  \code{I} | \code{isource} & von, nach & \code{wave} & Kurvenform (Abschnitt~\ref{sec:ref-wave}) \\
   & & \code{ac} & 0 \\
   & & \code{acphase} & 0 (Grad) \\
  \code{E} | \code{vcvs} & $+$, $-$, St$+$, St$-$ & \code{gain} & Pflicht \\
  \code{G} | \code{vccs} & von, nach, St$+$, St$-$ & \code{gm} | \code{gain} & Pflicht \\
  \code{F} | \code{cccs} & von, nach & \code{control} & Pflicht: Name einer V-Quelle oder Spule \\
   & & \code{gain} & Pflicht \\
  \code{H} | \code{ccvs} & $+$, $-$ & \code{control} & Pflicht \\
   & & \code{r} | \code{gain} & Pflicht (\ohm) \\
  \code{B} | \code{behavioral} & $+$, $-$ & genau eines von \code{i}, \code{v}, \code{q}, \code{phi} & Ausdruck: Strom, Spannung, Ladung, Fluss \\
  \code{K} | \code{coupling} & – & \code{inductors} & Pflicht: Liste von zwei Spulennamen \\
   & & \code{k} & Pflicht, $-1 < k < 1$ \\
\end{longtable}

\subsection*{Leitung, Schalter, ideale Diode, Teilschaltung}

\begin{longtable}{@{}P{2.7cm}P{3.0cm}P{3.6cm}P{4.5cm}@{}}
  \toprule
  \code{type} & \code{nodes} & Schlüssel & Standard / Bemerkung \\
  \midrule
  \endhead
  \bottomrule
  \endfoot
  \code{T} | \code{tline} & a$+$, a$-$, b$+$, b$-$ & \code{z0} & Pflicht (ohne \code{l}/\code{c}) \\
   & & \code{td} & Laufzeit; oder \code{f} mit \code{nl} \\
   & & \code{f}, \code{nl} & Frequenz; \code{nl} Standard 0,25 \\
   & & \code{l}, \code{c} & Beläge je Länge (statt \code{z0}, \code{td}) \\
   & & \code{len} & 1 \\
   & & \code{r}, \code{g} & 0, 0 (je Länge) \\
   & & \code{segments} & automatisch (höchstens 200) \\
  \code{S} | \code{switch} & a, b oder a, b, St$+$, St$-$ & \code{control} & Ausdruck; Pflicht bei zwei Knoten, verboten bei vier \\
   & & \code{vt} | \code{threshold} & 0 \\
   & & \code{vh} | \code{hysteresis} & 0 \\
   & & \code{ron}, \code{roff} & $10^{-3}$\,\ohm, $10^{9}$\,\ohm \\
   & & \code{initial} & \code{"off"}; oder \code{"on"} \\
  \code{ideal_diode} & Anode, Kathode & \code{vf} & 0 \\
   & & \code{ron}, \code{roff} & $10^{-3}$\,\ohm, $10^{9}$\,\ohm \\
  \code{X} | \code{subcircuit} & wie \code{ports} & \code{subcircuit} & Pflicht: Name \\
   & & \code{parameters} & \code{{"name": Zahl}}, überschreibt Standardwerte \\
\end{longtable}

\subsection*{Halbleiter}

\begin{longtable}{@{}P{3.2cm}P{3.4cm}P{7.2cm}@{}}
  \toprule
  \code{type} & \code{nodes} & Schlüssel \\
  \midrule
  \endhead
  \bottomrule
  \endfoot
  \code{D} | \code{diode} & Anode, Kathode & \code{model} und/oder Diodenparameter \\
  \code{Q} | \code{bjt} | \code{npn} | \code{pnp} & C, B, E & \code{model} und/oder BJT-Parameter; Polarität aus \code{type} oder Modell \\
  \code{J} | \code{jfet} | \code{njf} | \code{pjf} & D, G, S & \code{model} und/oder JFET-Parameter \\
  \code{M} | \code{mosfet} | \code{nmos} | \code{pmos} & D, G, S [, B] & \code{model}, \code{level}, \code{w}, \code{l}, MOS-Parameter; bei BSIM3 auch \code{ad}, \code{as}, \code{pd}, \code{ps}, \code{m} \\
\end{longtable}

Parameter am Bauteil überschreiben die des Modells. Ohne Bulk-Knoten liegt
der Bulk eines MOSFET an der Source.

% ---------------------------------------------------------------------------
\section{Kurvenformen (wave)}
\label{sec:ref-wave}

\begin{longtable}{@{}P{2.6cm}P{3.6cm}P{7.6cm}@{}}
  \toprule
  \code{type} & Schlüssel & Standard \\
  \midrule
  \endhead
  \bottomrule
  \endfoot
  \code{dc} & \code{value} & Pflicht \\
  \code{sin} | \code{sine} & \code{amplitude}, \code{freq} & Pflicht \\
   & \code{offset}, \code{delay}, \code{damping}, \code{phase} & 0 (\code{phase} in Grad, \code{damping} in 1/s) \\
  \code{pulse} & \code{v1}, \code{v2} & Pflicht \\
   & \code{delay} & 0 \\
   & \code{width}, \code{period} & unendlich (ohne \code{period}: Einzelpuls) \\
   & \code{rise}, \code{fall} & 0,1\,\% der Periode (sonst der Breite, sonst 1\,ns) \\
  \code{pwl} & \code{points} & Pflicht: \code{[[t, wert], ...]} \\
\end{longtable}

% ---------------------------------------------------------------------------
\section{Modelle}
\label{sec:ref-modelle}

Ein Modell ist ein Objekt mit \code{type} und Parametern, eine SPICE-Zeile
\code{.model NAME TYP(...)} oder eine Liste solcher Zeilen.

\begin{longtable}{@{}P{3.6cm}P{10.2cm}@{}}
  \toprule
  \code{type} & Parameter (Bedeutung und Standardwerte in Kapitel~\ref{kap:format-modelle}) \\
  \midrule
  \endhead
  \bottomrule
  \endfoot
  \code{diode} | \code{D} & \code{is n rs cj0 vj m fc tt bv ibv eg xti kf af} \\
  \code{npn}, \code{pnp} & \code{is bf br nf nr vaf var ikf ikr ise ne isc nc rb rc re cje vje mje cjc vjc mjc tf tr fc eg xti xtb kf af} \\
  \code{njf}, \code{pjf} & \code{vto beta lambda is cgs cgd pb fc kf af} \\
  \code{nmos}, \code{pmos}, Level 1 & \code{vto kp gamma phi lambda w l tox cgso cgdo cgbo cbd cbs pb mj fc is tcv kf af} \\
  \quad \code{"level": "ekv"} & zusätzlich \code{theta ucrit} \\
  \quad \code{"level": 3} & zusätzlich \code{theta eta vmax nfs kappa xj ld delta u0 nsub nss tpg} \\
  \quad \code{"level": "bsim3"} (SPICE \code{LEVEL=8}, \code{49}, \code{53}) & siehe unten \\
\end{longtable}

\paragraph{BSIM3.}
Ausgewertet werden:
\code{a0 a1 a2 acde ags alpha0 alpha1 at b0 b1 beta0 binunit capmod cdsc
cdscb cdscd cf cgbo cgdl cgdo cgsl cgso cit cj cjsw cjswg ckappa clc cle
delta dlc drout dsub dvt0 dvt0w dvt1 dvt1w dvt2 dvt2w dwb dwc dwg eta0 etab
gamma1 gamma2 js jsw k1 k2 k3 k3b keta kf af kt1 kt1l kt2 lint ll llc lln
lw lwc lwl lwlc lwn mj mjsw mjswg mobmod moin nfactor ngate nj nlx noff
npeak nsub pb pbsw pbswg pclm pdibl1 pdibl2 pdiblb prt prwb prwg pscbe1
pscbe2 pvag rdsw tnom tox toxm u0 ua ua1 ub ub1 uc uc1 ute vbm vbx vfb vfbcv
voff voffcv vsat vth0 w0 wint wl wlc wln wr ww wwc wwl wwlc wwn xj xl xpart
xt xti xw}; am Bauteil \code{w l ad as pd ps m}. Die größenabhängigen
Parameter gibt es zusätzlich mit den Vorsilben \code{l}, \code{w} und
\code{p} (z.\,B. \code{lvth0}). Synonyme: \code{vtho} = \code{vth0},
\code{nch} = \code{npeak}, \code{pdiblc1}/\code{pdiblc2}/\code{pdiblcb} =
\code{pdibl1}/\code{pdibl2}/\code{pdiblb}, \code{cjo} = \code{cj}.
Ohne Wirkung, aber ohne Warnung angenommen: \code{noimod version paramchk
lmin lmax wmin wmax tcj tcjsw tcjswg tpb tpbsw tpbswg elm rsh hdif nqsmod
acnqsmod noia noib noic em ef lref wref acm ijth wmlt lmlt nrd nrs geo level
lelm welm pelm tox1 dtoxcv}.

% ---------------------------------------------------------------------------
\section{Teilschaltungen}
\label{sec:ref-teil}

\begin{longtable}{@{}P{3.0cm}P{4.3cm}P{6.5cm}@{}}
  \toprule
  Schlüssel & Wert & Bedeutung \\
  \midrule
  \endhead
  \bottomrule
  \endfoot
  \code{ports} & Liste von Knotennamen, Pflicht & Anschlüsse nach außen \\
  \code{parameters} & \code{{"name": Zahl}} & Standardwerte der Parameter \\
  \code{components} & Liste & Bauteile, auch Instanzen (bis 20 Ebenen) \\
  \code{variables}, \code{equations}, \code{events} & wie oben & eigene Gleichungen und Ereignisse \\
  \code{schematic} & Objekt & Plan des Editors \\
\end{longtable}

Statt eines Objekts darf der SPICE-Text einer \code{.subckt} stehen
(Zeichenkette oder Liste von Zeilen, Abschnitt~\ref{sec:spice-subckt}).
Namen von außen: \code{<instanz>.<name>}, geschachtelt
\code{<instanz>.<instanz>.<name>}.

% ---------------------------------------------------------------------------
\section{Gleichungen, Ereignisse, Tabellen}
\label{sec:ref-gleichungen}

\subsection*{variables}

\begin{longtable}{@{}P{4.6cm}P{9.2cm}@{}}
  \toprule
  Form & Bedeutung \\
  \midrule
  \endhead
  \bottomrule
  \endfoot
  \code{"x": Zahl} & Anfangswert (Zustand) bzw. Startwert \\
  \code{"x": null} & Zustand startet im Gleichgewicht \\
  \code{"x": {"init": Zahl}} & Anfangswert \\
  \code{"x": {"guess": Zahl}} & nur Startwert für Newton \\
  \code{"x": {"discrete": Zahl}} & diskrete Variable mit Startwert \\
  \code{"abstol": Zahl} & zusätzlich in einem Objekt: eigene absolute Toleranz \\
  \code{"T[n]": ...} & Vektor \code{T[1]} … \code{T[n]}; Werte als Zahl oder Liste je Element \\
\end{longtable}

Von \code{init}, \code{guess} und \code{discrete} darf höchstens eines
stehen.

\subsection*{equations}

\begin{longtable}{@{}P{6.0cm}P{7.8cm}@{}}
  \toprule
  Form & Bedeutung \\
  \midrule
  \endhead
  \bottomrule
  \endfoot
  \code{"links = rechts"} & eine Gleichung; \code{der(...)} für Zeitableitungen \\
  \code{{"for": "i", "from": a, "to": b, "equation": "..."}} & eine Gleichung je Indexwert \\
  \code{{"for": "i", "from": a, "to": b, "equations": [...]}} & mehrere Einträge je Indexwert, auch Schleifen \\
\end{longtable}

\subsection*{events}

\begin{longtable}{@{}P{3.0cm}P{4.3cm}P{6.5cm}@{}}
  \toprule
  Schlüssel & Wert & Bedeutung \\
  \midrule
  \endhead
  \bottomrule
  \endfoot
  \code{when} & Ausdruck mit Vergleich & feuert, wenn die Bedingung wahr wird \\
  \code{every} & Zahl $>0$ & feuert periodisch \\
  \code{start} & Zahl, Standard: eine Periode & erster Zeitpunkt von \code{every} \\
  \code{at} & Zahl oder Liste & feuert zu diesen Zeiten \\
  \code{set} & \code{{"ziel": "Ausdruck"}} & Zuweisungen an Zustände, diskrete Variablen, \code{v(knoten)}, \code{i(spule)} \\
\end{longtable}

Genau einer der Auslöser \code{when}, \code{every}, \code{at}.

\subsection*{tables}
\label{sec:ref-tabellen}

\begin{longtable}{@{}P{3.0cm}P{4.3cm}P{6.5cm}@{}}
  \toprule
  Schlüssel & Standard & Bedeutung \\
  \midrule
  \endhead
  \bottomrule
  \endfoot
  \code{x}, \code{y} & – & Stützstellen als zwei Listen \\
  \code{points} & – & Stützstellen als Liste von \code{[x, y]} \\
  \code{file} & – & Textdatei mit Spalten, relativ zur Eingabedatei \\
  \code{xcol}, \code{ycol} & 0, 1 & Spalten der Datei (ab 0) \\
  \code{interp} & \code{"linear"} & oder \code{"cubic"} (monoton, PCHIP) \\
  \code{extrap} & \code{"hold"} & oder \code{"linear"} \\
\end{longtable}

\subsection*{functions}

\code{"name(a, b, ...)": "Ausdruck"} – Aufruf in Ausdrücken als
\code{name(...)}; keine Rekursion, keine eingebauten Namen.

% ---------------------------------------------------------------------------
\section{Analysen}
\label{sec:ref-analysen}

Alle Analysen haben \code{type} (Pflicht) und \code{name}. Dazu kommen
\code{outputs} bei \code{op}, \code{dc}, \code{ac}, \code{tran} und
\code{pss}; \code{measure} bei \code{dc}, \code{ac}, \code{tran} und
\code{pss}; \code{save} (Standard \code{true}) bei \code{dc}, \code{ac},
\code{tran}, \code{pss} und \code{noise}; \code{param_sweep} oder
\code{monte_carlo} bei allen außer \code{sens} und \code{fit}.

\begin{longtable}{@{}P{1.8cm}P{3.6cm}P{3.4cm}P{5.0cm}@{}}
  \toprule
  \code{type} & Schlüssel & Standard & Bedeutung \\
  \midrule
  \endhead
  \bottomrule
  \endfoot
  \code{op} & – & & Arbeitspunkt \\
  \addlinespace[0.8ex]
  \code{dc} & \code{source} & Pflicht & V- oder I-Quelle \\
   & \code{start}, \code{stop}, \code{step} & Pflicht & Bereich \\
  \addlinespace[0.8ex]
  \code{ac} & \code{fstart}, \code{fstop} & Pflicht & Frequenzbereich \\
   & \code{sweep} & \code{"dec"} & \code{"dec"}, \code{"oct"}, \code{"lin"} \\
   & \code{points} & 10 & je Dekade/Oktave bzw. insgesamt \\
  \addlinespace[0.8ex]
  \code{tran} & \code{tstop} & Pflicht & Endzeit \\
   & \code{tstep} & \code{tstop/50} & größter Schritt (Auflösung) \\
   & \code{hmax}, \code{hmin}, \code{hinit} & aus \code{tstep}, 0, 0 & Schrittweiten \\
   & \code{method} & \code{"bdf"} & \code{"bdf"}, \code{"gear2"}, \code{"trap"}, \code{"be"} \\
   & \code{max_order} & 5 & höchste BDF-Ordnung \\
   & \code{uic} & \code{false} & ohne Arbeitspunkt starten \\
   & \code{output_step} & – & festes Ausgaberaster \\
   & \code{max_events} & 100000 & Höchstzahl der Ereignisse \\
   & \code{fourier} & – & Abschnitt~\ref{sec:ref-fourier} \\
  \addlinespace[0.8ex]
  \code{noise} & \code{output} & Pflicht & \code{v(k)}, \code{v(a,b)}, \code{i(b)} \\
   & \code{input} & – & Quelle für das eingangsbezogene Rauschen \\
   & \code{fstart}, \code{fstop}, \code{sweep}, \code{points} & wie \code{ac} & \\
  \addlinespace[0.8ex]
  \code{pss} & \code{period} & Pflicht & Periode der Anregung \\
   & \code{periods} & 1 & ausgegebene Perioden \\
   & \code{tol} & $10^{-4}$ & relative Abweichung \\
   & \code{max_iterations} & 60 & \\
   & \code{tstab} & 0 & Vorlauf vor der Suche \\
   & alle Schlüssel von \code{tran} außer \code{tstop} & & \\
  \addlinespace[0.8ex]
  \code{sens} & \code{analysis} & Pflicht & eingebettete Analyse \\
   & \code{parameters} & alle & Liste von Parameternamen \\
   & \code{delta} & $10^{-3}$ & relative Änderung \\
  \addlinespace[0.8ex]
  \code{fit} & \code{analysis} & Pflicht & \code{dc}, \code{ac} oder \code{tran} \\
   & \code{output} & Pflicht & verglichene Größe \\
   & \code{data} & Pflicht & \code{x}/\code{y}, \code{points} oder \code{file}/\code{xcol}/\code{ycol} \\
   & \code{parameters} & Pflicht & \code{{"p": {"start", "min", "max", "log"}}} \\
   & \code{residual} & \code{"lin"} & \code{"lin"}, \code{"rel"}, \code{"log"} \\
   & \code{max_iterations} & 50 & \\
   & \code{tol} & $10^{-8}$ & Abbruch der Iteration \\
\end{longtable}

\subsection*{param\_sweep}

Ein Objekt oder eine Liste von Objekten (verschachtelt, das erste außen):

\begin{longtable}{@{}P{3.0cm}P{3.4cm}P{7.4cm}@{}}
  \toprule
  Schlüssel & Standard & Bedeutung \\
  \midrule
  \endhead
  \bottomrule
  \endfoot
  \code{param} & Pflicht & Name aus \code{parameters} \\
  \code{values} & – & Liste der Werte, oder: \\
  \code{start}, \code{stop}, \code{points} & – & Bereich mit mindestens 2 Punkten \\
  \code{scale} & \code{"lin"} & oder \code{"log"} \\
\end{longtable}

\subsection*{monte\_carlo}

\begin{longtable}{@{}P{3.0cm}P{3.4cm}P{7.4cm}@{}}
  \toprule
  Schlüssel & Standard & Bedeutung \\
  \midrule
  \endhead
  \bottomrule
  \endfoot
  \code{runs} & 100 & Zahl der Läufe \\
  \code{seed} & 1 & Startwert des Zufallsgenerators \\
  \code{vary} & Pflicht & \code{{"param": {...}}} je gestreutem Parameter: \\
  \quad\code{tol} & – & relative Toleranz, oder: \\
  \quad\code{sigma} & – & absolute Standardabweichung \\
  \quad\code{dist} & \code{"uniform"} (mit \code{sigma}: \code{"gauss"}) & \code{"uniform"} oder \code{"gauss"} \\
\end{longtable}

Monte Carlo braucht \code{measure}; die Kurven der Läufe werden nur mit
\code{"save": true} gespeichert.

\subsection*{measure}
\label{sec:ref-measure}

\code{{"name": {"funktion": "Ausdruck", option: wert, ...}, ...}} mit genau
einer Funktion:

\begin{longtable}{@{}P{3.6cm}P{10.2cm}@{}}
  \toprule
  Funktion & Optionen \\
  \midrule
  \endhead
  \bottomrule
  \endfoot
  \code{min max pp argmin argmax avg rms integ final} & \code{from}, \code{to} \\
  \code{at} & \code{x} (Pflicht), \code{from}, \code{to} \\
  \code{cross} & \code{level} (Pflicht), \code{n} (1), \code{edge} (\code{"rise"}, \code{"fall"}, \code{"both"}), \code{from}, \code{to} \\
  \code{period freq duty} & \code{level} (Mitte), \code{from}, \code{to} \\
  \code{rise fall} & \code{low}, \code{high} (10\,\%, 90\,\% des Hubs), \code{from}, \code{to} \\
  \code{f3db f3db_low} & \code{from}, \code{to} \\
  \code{thd harmonic} & \code{fundamental}, \code{n}, \code{from}, \code{to} \\
\end{longtable}

\subsection*{fourier}
\label{sec:ref-fourier}

\begin{longtable}{@{}P{3.0cm}P{3.4cm}P{7.4cm}@{}}
  \toprule
  Schlüssel & Standard & Bedeutung \\
  \midrule
  \endhead
  \bottomrule
  \endfoot
  \code{freq} & Pflicht & Grundfrequenz \\
  \code{harmonics} & 9 & Zahl der Harmonischen \\
  \code{periods} & 1 & ausgewertete Perioden am Ende \\
  \code{outputs} & die Ausgänge & ausgewertete Größen \\
  \code{from} & 0 & frühester Zeitpunkt \\
\end{longtable}

% ---------------------------------------------------------------------------
\section{Optionen}
\label{sec:ref-optionen}

\begin{longtable}{@{}P{3.4cm}P{2.6cm}P{7.8cm}@{}}
  \toprule
  Schlüssel & Standard & Bedeutung \\
  \midrule
  \endhead
  \bottomrule
  \endfoot
  \code{reltol} & $10^{-4}$ & relative Toleranz \\
  \code{vntol} & $10^{-6}$ & absolute Toleranz für Spannungen (V) \\
  \code{abstol} & $10^{-12}$ & absolute Toleranz für Ströme (A) \\
  \code{vartol} & $10^{-9}$ & absolute Toleranz für Variablen \\
  \code{lte_floor} & 0,1 & Bezugsgröße des Zeitschrittfehlers \\
  \code{trtol} & 1 & Faktor auf den zulässigen Zeitschrittfehler \\
  \code{gmin} & $10^{-12}$ & Leitwert parallel zu pn-Übergängen (S) \\
  \code{temp} & 27 & Temperatur der Schaltung (°C) \\
  \code{tnom} & 27 & Temperatur der Modellparameter (°C) \\
  \code{maxiter_dc} & 100 & Newton-Iterationen im Arbeitspunkt \\
  \code{maxiter_tran} & 20 & Newton-Iterationen je Zeitschritt \\
  \code{reuse_jacobian} & \code{true} & zerlegte Matrix im Schritt weiterverwenden \\
  \code{index_reduction} & \code{true} & Index über 2 automatisch reduzieren \\
\end{longtable}

% ---------------------------------------------------------------------------
\section{Ausdrücke}

\begin{longtable}{@{}P{3.6cm}P{10.2cm}@{}}
  \toprule
  Element & Schreibweise \\
  \midrule
  \endhead
  \bottomrule
  \endfoot
  Operatoren & \code{+ - * /}, \code{^} | \code{**}, \code{< <= > >= == !=}, \code{&& || !}, Klammern \\
  Funktionen (1 Arg.) & \code{sin cos tan asin acos atan sinh cosh tanh exp log ln log10 sqrt abs sign} \\
  Funktionen (2 Arg.) & \code{atan2(y, x) min(a, b) max(a, b) pow(a, b)} \\
  Fallunterscheidung & \code{if(bedingung, dann, sonst)} \\
  Besondere & \code{der(x)} (nur Gleichungen), \code{delay(x, td)}, \code{sum(i, von, bis, ausdruck)} \\
  Namen & \code{t} | \code{time}, \code{pi}, Parameter, Variablen, \code{T[i]}, Tabellen, eigene Funktionen \\
  Schaltung & \code{v(knoten)}, \code{v(a, b)}, \code{i(bauteil)}, mit \code{instanz.} für Teilschaltungen \\
  Ausgänge & zusätzlich \code{T[*]} für alle Elemente eines Vektors \\
  SI-Präfixe & \code{f p n u}, \texttt{µ}, \code{m k K M meg G T} (Abschnitt~\ref{sec:zahlen}) \\
\end{longtable}

% ---------------------------------------------------------------------------
\section{Befehle des Editors}

\subsection*{Menüs}

\begin{longtable}{@{}P{2.6cm}P{6.6cm}P{4.6cm}@{}}
  \toprule
  Menü & Befehl & Kürzel \\
  \midrule
  \endhead
  \bottomrule
  \endfoot
  Datei & Neu & \tasten{Strg}{N} \\
   & Öffnen \dots & \tasten{Strg}{O} \\
   & Zuletzt geöffnet \pfad\ \textit{Datei}, Liste leeren & \\
   & Beispiele \pfad\ \textit{Beispiel} & \\
   & Sichern & \tasten{Strg}{S} \\
   & Sichern unter \dots & \tasten{Umschalt+Strg}{S} \\
   & Beenden & \tasten{Strg}{Q} \\
  Bearbeiten & Widerrufen & \tasten{Strg}{Z} \\
   & Wiederholen & \tasten{Strg}{Y} \\
   & Ausschneiden, Kopieren, Einfügen & \tasten{Strg}{X}, \tasten{Strg}{C}, \tasten{Strg}{V} \\
   & Alles auswählen & \tasten{Strg}{A} \\
   & Kommentar ein/aus & \tasten{Strg}{/} \\
   & Auswahl in Teilschaltung umwandeln \dots & \tasten{Strg}{G} \\
   & Modellbibliothek \dots & \tasten{Strg}{L} \\
  Simulation & Starten & \tasten{Strg}{R} \\
   & Stoppen & \tasten{Strg}{.} \\
  Ansicht & Schaltplan, Text & \\
   & Schaltplan einpassen & (\taste{F} im Plan) \\
   & Bauteile, Eigenschaften, Ergebnisse, Protokoll, Probleme & \\
   & Kurven: alles zeigen & \tasten{Strg}{0} \\
  Hilfe & Über ESIM Editor & \\
\end{longtable}

\subsection*{Kontextmenü im Plan}

\begin{longtable}{@{}P{6.6cm}P{2.6cm}P{4.6cm}@{}}
  \toprule
  Befehl & Kürzel & erscheint \\
  \midrule
  \endhead
  \bottomrule
  \endfoot
  Teilschaltung „…“ öffnen / SPICE-Text zeigen & Doppelklick & auf einer Instanz \\
  Eigenschaften von \textit{Bauteil} & Doppelklick & auf einem Bauteil \\
  Strom i(\textit{Bauteil}) zeigen & \tasten{Alt}{Klick} & auf einem Bauteil \\
  Spannung v(\textit{Netz}) zeigen & \tasten{Alt}{Klick} & auf einem Netz \\
  Umbenennen \dots & Doppelklick & auf einem Label \\
  Drehen & \taste{R} & immer \\
  Spiegeln & \taste{X} & immer \\
  Verdoppeln & \tasten{Strg}{D} & immer \\
  Löschen & \taste{Entf} & immer \\
  In Teilschaltung umwandeln \dots & \tasten{Strg}{G} & immer \\
  Leitung zeichnen & \taste{W} & immer \\
  Masse & \taste{G} & immer \\
  Netzlabel \dots & \taste{N} & immer \\
  Port \dots & \taste{P} & in Teilschaltungen \\
  Bauteil \dots & \taste{A} & immer \\
  Alles zeigen & \taste{F} & immer \\
  Zur Hauptschaltung & \tasten{Strg}{$\uparrow$} & in Teilschaltungen \\
\end{longtable}

\subsection*{Einträge der Bauteilliste}

Leitung zeichnen, Masse, Netzlabel, Port der Teilschaltung (nur in
Teilschaltungen), Widerstand, Kondensator, Spule, Spannungsquelle,
Sinusquelle, Pulsquelle, Stromquelle, Diode, Ideale Diode, NPN-Transistor,
PNP-Transistor, N-MOSFET, P-MOSFET, N-JFET, P-JFET, Spannungsgesteuerte
Spannungsquelle, Spannungsgesteuerte Stromquelle, Stromgesteuerte
Stromquelle, Stromgesteuerte Spannungsquelle, Verhaltensquelle (Spannung),
Verhaltensquelle (Strom), Schalter (spannungsgesteuert), Leitung, Kopplung
zweier Spulen, je eine Teilschaltung der Datei, Modellbibliothek \dots, Neue
Teilschaltung \dots; beim Suchen zusätzlich alle Modelle der Bibliothek.

\subsection*{Kurven und Histogramme}

Rechtsklick: \menue{Alles zeigen}, \menue{Als Bild speichern \dots};
Auswahl \menue{Darstellung}: \menue{Histogramm}, \menue{Streuung über
\textit{p}} (Monte Carlo), \menue{über \textit{p}} (Sweeps).

\subsection*{Dialog Modellbibliothek}

Suchfeld, Art (Alle Arten, Dioden, Bipolartransistoren, MOSFETs, JFETs,
Teilschaltungen (SPICE)), \menue{Bibliotheken \dots} (\menue{Datei \dots},
\menue{Ordner \dots}, \menue{Entfernen}), \menue{Im Plan platzieren},
\menue{In die Datei übernehmen}, \menue{Verwenden}.
